\scanpage{97}
comme
\[
\left\{\begin{aligned}
&T_{\Sigma_I,t}\simeq O_t(2)(\underline\omega)\\
&\text{avec }\ O_t(2)\overset{\mathrm{def}}{\simeq}O_t(1)^{\otimes2},\qquad
O_t(1)\simeq V_I/F_t\simeq F_t^\vee(\underline\omega).
\end{aligned}\right. \tag*{(62)}
\]
de sorte que $\unsure{\nu_E|v_E}$ s'interprète comme un isom
\[
\begin{aligned}
\nu:\underline t\deleted{\unsure{_E}}^{\otimes2}&\simeq
\overset{\deleted{\unsure{\vee}}}{F}_t^{\otimes(-2)}(\underline\omega)
\qquad\text{ou aussi comme}\\[-2pt]
&&\check\nu:\underline t^{\otimes-2}\simeq F_t^{\otimes2}(\underline\omega)
\end{aligned}\tag*{(63)}
\]
(NB $\underline\omega=\underline\omega_{\underline I}$, ens des deux ordres circulaires sur $\underline I$). Les données de base sont donc pour nous
\[
\begin{aligned}
[\,I\;(&\mapsto\underline\omega=\underline\omega_I\deleted{\unsure{)}},\deleted{\unsure{F}}\text{ et }V_I=V_I(k)),\\
&F_t\hookrightarrow V_I,\ \underline t,\ \nu:t^{\otimes2}\simeq F_t^{\otimes(-2)}(\underline\omega)\,].
\end{aligned}
\]
Ceci posé, la donnée d'un épinglage de Legendre s'explicite comme un isom
\[
\kappa:\underline t\xrightarrow{\sim}\underline O_t(1)^{\unsure{\sharp|\xi|\natural}} \tag*{(64)}
\]
i.e. comme une famille d'isomorphismes
\[
\left\{\begin{aligned}
&(\kappa_i)_{i\in\mu_4^*},\qquad
\kappa_i:\underline t\xrightarrow{\sim}\underline O_t(1)\simeq F_t^\vee(\underline\omega)\\
&\kappa_{i'}=i\,\kappa_i\qquad\text{si }i'=1/i\text{ est l'autre élément}\\
&\hphantom{\kappa_{i'}=i\,\kappa_i\qquad}\text{de }\mu_4^\times
\end{aligned}\right. \tag*{(65)}
\]
% slanted
\margin{$\mu_4^\times=\mu_4^*(k)$, éléments d'ordre $4$ de $k^*$}
ou encore, d'une famille
\[
\left\{\begin{aligned}
&(\check\kappa_i)_{i\in\mu_4^\times},\qquad \check\kappa_i:\underline t^\vee\simeq F_t(\underline\omega)\\
&\check\kappa_{i'}=\frac1i\,\check\kappa_i\;,
\end{aligned}\right. \tag*{(66)}
\]
satisfaisant à une certaine condition :
\[
\left\{\begin{aligned}
&\kappa_i^{\otimes2}:t^{\otimes2}\xrightarrow{\sim}\overset{\deleted{\unsure{\vee|\wedge}}}{F}_t^{\otimes(-2)}
\quad\text{est de la forme}\\
&\kappa_i^{\otimes2}(\unsure{\delta|\vartheta}^{\otimes2})=\nu(\delta^{\otimes2})\mathbin{\wedge}\omega_i
\end{aligned}\right. \tag*{(67)}
\]
où $\omega_i\in\underline\omega$ est un élément convenable de $\underline\omega$, bien déterminé bien sûr par $\kappa,i$, en fait par $\kappa_i$.
% archive sheet number: 96
